\documentclass[12pt]{article}

\usepackage[utf8]{inputenc}
\usepackage[bulgarian]{babel}
\usepackage{amssymb}
\usepackage{amsmath}
\usepackage{hyperref}
\usepackage{booktabs}
\usepackage{listings}
\usepackage{xcolor}
\usepackage[top=1.3in, bottom=1.5in, left=1.3in, right=1.3in]{geometry}

\definecolor{codebg}{rgb}{0.95,0.95,0.95}
\lstset{
  backgroundcolor=\color{codebg},
  basicstyle=\ttfamily\small,
  breaklines=true,
  frame=single,
}

\begin{document}
    \begin{center}
        \LARGE{\textbf{Тема: XML2Emmet — разгръщане в AWS с EC2, RDS и S3}}

        \bigskip
        \Large{Предмет: Приложно-програмни интерфейси за работа с облачни архитектури с Amazon Web Services (AWS)}

        \medskip
        \Large{Изготвил: Денис Мустафа Каим, фн: 3MI0600294, имейл: deniskaim86@gmail.com}

        \medskip
        \Large{Лектор: проф. д-р Милен Петров, година: 2026}

        \bigskip
    \end{center}

    \tableofcontents
    \bigskip
    \newpage


\section{Условие}

\noindent Задачата е да се разгърне съществуващ уеб проект от курса по Уеб технологии в Amazon Web Services, като се добави нова функционалност, използваща специфични AWS услуги. Всеки студент разширява проекта самостоятелно, като избира услуги, различни от тези на останалите членове на групата.

\medskip

\noindent Проектът \textbf{xml2emmet} е уеб приложение за трансформация на \textbf{XML} към \textbf{Emmet} нотация, разработен в курса по Уеб технологии. За AWS разширението е добавена функционалност за export на история в Amazon S3, а приложението е хоствано на Amazon EC2 с база данни в Amazon RDS.


\section{Въведение}

xml2emmet позволява на потребителите да трансформират XML и HTML структури в Emmet нотация. Приложението е изградено с PHP 8.2, REST API и SPA frontend.

\noindent За разгръщането в AWS са използвани следните услуги:

\begin{itemize}
    \item \textbf{Amazon EC2} — хостване на уеб сървъра.

    \item \textbf{Amazon RDS MySQL} — съхранение на потребителски данни и история.

    \item \textbf{Amazon S3} — export на история като JSON файл с presigned URL.

    \item \textbf{SSM Parameter Store} — сигурно съхранение на credentials.

    \item \textbf{CloudWatch Logs} — логове на приложението.

    \item \textbf{CloudFormation} — инфраструктурата като код.
\end{itemize}


\section{Теория}

\noindent \textbf{Amazon EC2} (Elastic Compute Cloud) е услуга за виртуални машини в облака. Позволява стартиране на сървъри с различни операционни системи и конфигурации. В проекта се използва инстанция тип t3.micro с Amazon Linux 2023.

\medskip

\noindent \textbf{Amazon RDS} (Relational Database Service) управлява релационни бази данни в облака. Осигурява автоматични backup-и, мониторинг и висока наличност. В проекта се използва MySQL 8.0 на db.t3.micro инстанция.

\medskip

\noindent \textbf{Amazon S3} (Simple Storage Service) е обектно хранилище с висока надеждност и мащабируемост. Поддържа presigned URL-и — временни линкове за достъп до обекти без публичен достъп до bucket-а.

\medskip

\noindent \textbf{SSM Parameter Store} позволява сигурно съхранение на конфигурационни стойности и тайни. EC2 инстанцията ги чете при стартиране чрез AWS CLI.

\medskip

\noindent \textbf{AWS CloudFormation} описва цялата инфраструктура като YAML шаблон. Един файл създава VPC, подмрежи, security groups, RDS, S3, SSM параметри и EC2.


\section{Използвани технологии}

\begin{itemize}
    \item \textbf{Amazon EC2} - t3.micro, Amazon Linux 2023, Apache 2.4, PHP 8.2.

    \item \textbf{Amazon RDS MySQL} - версия 8.0, db.t3.micro, 20 GB gp2.

    \item \textbf{Amazon S3} - обектно хранилище за JSON export файлове.

    \item \textbf{Amazon VPC} - CIDR 10.0.0.0/16, публична и частна подмрежа.

    \item \textbf{SSM Parameter Store} - съхранение на DB credentials и S3 bucket.

    \item \textbf{CloudWatch Logs} - лог група \texttt{/ec2/xml2emmet}, ретенция 7 дни.

    \item \textbf{AWS CloudFormation} - IaC шаблон за цялата инфраструктура.

    \item \textbf{aws/aws-sdk-php} - версия \^{}3.318, PHP библиотека за AWS услуги.
\end{itemize}


\section{Инсталация и настройки}

\noindent\textbf{Стъпка 1.} Отвори AWS CloudShell от конзолата на AWS Academy Learner Lab и клонирай репото:

\begin{lstlisting}[language=bash]
git clone -b feat/aws-ec2-rds-s3 \
  https://github.com/peterdobrev/xml2emmet.git
cd xml2emmet
\end{lstlisting}

\medskip

\noindent\textbf{Стъпка 2.} Разгърни CloudFormation стека с една команда. Стекът създава цялата инфраструктура автоматично:

\begin{lstlisting}[language=bash]
aws cloudformation deploy \
  --template-file deploy/cloudformation.yml \
  --stack-name xml2emmet \
  --parameter-overrides \
      KeyPairName=vockey \
      DBMasterPassword=<парола> \
      GithubRepoUrl=https://github.com/peterdobrev/xml2emmet.git \
      AppBranch=feat/aws-ec2-rds-s3 \
  --region us-east-1
\end{lstlisting}

\medskip

\noindent\textbf{Стъпка 3.} Изчакай ~12 минути (RDS отнема най-много). След \texttt{CREATE\_COMPLETE} вземи URL-а:

\begin{lstlisting}[language=bash]
aws cloudformation describe-stacks \
  --stack-name xml2emmet \
  --query 'Stacks[0].Outputs[?OutputKey==`AppURL`].OutputValue' \
  --output text --region us-east-1
\end{lstlisting}



\section{Кратко ръководство за потребителя}

Приложението е достъпно на адреса, върнат от CloudFormation. Потребителят отваря \texttt{/app.html} в браузъра, регистрира се и може да:

\begin{itemize}
    \item трансформира XML/HTML към Emmet нотация,
    \item преглежда история на трансформациите,
    \item експортира историята си като JSON файл в S3.
\end{itemize}


\noindent За S3 export потребителят извиква \texttt{POST /api/history/export}. Отговорът съдържа presigned URL, валиден 1 час:

\begin{lstlisting}[language=bash]
curl -s -X POST http://<APP_URL>/api/history/export \
  -H "Cookie: xml2emmet_sid=<session>"
\end{lstlisting}



\section{Примерни данни}

Примерна трансформация:

\begin{table}[h]
\caption{Вход и изход на трансформацията}
\center
\begin{tabular}{ |p{6cm}|p{6cm}| }
 \hline
 \textbf{XML вход} & \textbf{Emmet изход} \\ \hline
 \texttt{<ul><li>Item</li></ul>} & \texttt{ul>li\{Item\}} \\ \hline
 \texttt{<div class="box"><p>Text</p></div>} & \texttt{div.box>p\{Text\}} \\ \hline
\end{tabular}
\end{table}

\noindent Примерен S3 export файл (\texttt{exports/user-1/2026-06-09T070004Z.json}):

\begin{lstlisting}
[
  {
    "id": 1,
    "input": "<ul><li>Item</li></ul>",
    "output": "ul>li{Item}",
    "created_at": "2026-06-09 07:00:00"
  }
]
\end{lstlisting}


\section{Описание на програмния код}

\subsection{HistoryExporter.php}

Файлът \texttt{src/Aws/HistoryExporter.php} използва AWS SDK за PHP за качване на JSON в S3 и генериране на presigned URL:

\begin{lstlisting}[language=php]
$s3->putObject([
    'Bucket' => $this->bucket,
    'Key'    => "exports/user-{$userId}/{$timestamp}.json",
    'Body'   => json_encode($rows),
]);
$cmd = $s3->getCommand('GetObject', [
    'Bucket' => $this->bucket,
    'Key'    => $key,
]);
$url = $s3->createPresignedRequest($cmd, '+1 hour')->getUri();
\end{lstlisting}

\subsection{CloudFormation шаблон}

Файлът \texttt{deploy/cloudformation.yml} описва цялата инфраструктура. EC2 инстанцията съдържа UserData скрипт, който при първо стартиране инсталира Apache, PHP, клонира репото, чете credentials от SSM и стартира приложението.


\section{Приноси на студента, ограничения и възможности за бъдещо развитие}

В рамките на проекта са реализирани: разгръщане на приложението в AWS чрез CloudFormation, интеграция с Amazon S3 за export на история с presigned URL, и използване на SSM Parameter Store за сигурно управление на credentials.

\medskip

\noindent Ограничение е, че приложението е достъпно само по HTTP (без SSL сертификат). Като бъдещо развитие може да се добави Amazon CloudFront с ACM сертификат за HTTPS, или Auto Scaling Group за хоризонтално мащабиране.


\section{Какво научих}

В хода на проекта се запознах практически с разгръщането на уеб приложение в AWS — от конфигурацията на VPC и Security Groups до автоматизирането на целия процес с CloudFormation. Научих как EC2 инстанция може да получава credentials сигурно чрез SSM Parameter Store без hardcoded пароли в кода. Разбрах как работят presigned URL-ите в S3 и как Instance Profile дава права на EC2 да достъпва AWS услуги без допълнителна конфигурация.


\section{Списък с фигури и таблици}

\listoftables


\section{Използвани източници}

\noindent\href{https://docs.aws.amazon.com/cloudformation/}{[1] AWS CloudFormation Documentation}

\medskip

\noindent\href{https://docs.aws.amazon.com/AWSEC2/latest/UserGuide/}{[2] Amazon EC2 User Guide}

\medskip

\noindent\href{https://docs.aws.amazon.com/AmazonRDS/latest/UserGuide/}{[3] Amazon RDS User Guide}

\medskip

\noindent\href{https://docs.aws.amazon.com/AmazonS3/latest/userguide/}{[4] Amazon S3 User Guide}

\medskip

\noindent\href{https://docs.aws.amazon.com/systems-manager/latest/userguide/systems-manager-parameter-store.html}{[5] AWS Systems Manager Parameter Store}

\medskip

\noindent\href{https://github.com/aws/aws-sdk-php}{[6] AWS SDK for PHP — GitHub}

\end{document}
